
\begin{itemize}
\item El sistema evaluado presenta heterogeneidad estadísticamente significativa entre bibliotecas en iluminación, ruido y conectividad; no es metodológicamente válido tratarlas como unidades equivalentes.
\item El ruido es la dimensión más crítica para el confort de estudio, con incumplimientos recurrentes y asociación positiva con ocupación.
\item La conectividad muestra mejor desempeño relativo, pero con asimetrías entre sedes y eventos atípicos que afectan continuidad de servicio.
\item El índice ponderado propuesto mejora la sensibilidad diagnóstica frente al índice simple, al reflejar mayor peso de variables de confort ambiental para actividades de lectura y estudio.
\item Se recomienda institucionalizar monitoreo periódico y tablero de control por sede, integrando indicadores de tendencia central, dispersión, cumplimiento e incidencia de outliers.
\end{itemize}
